%----------------------------------------
\begin{frame}
  \frametitle{Conclusion} \small

\begin{enumerate}
\item A task network is a {\bf logical formula}. Scheduling is $\exists\exists$; verification is $\forall\forall$.
\pause
\item The interesting questions are the {\bf other quantifier orders}. $\forall\exists$ closes the
      vacuity trap that $\forall\forall$ leaves open.
\pause
\item The zone encoding is linear except for {\bf rate $\times$ duration}. That one product sets
      the scaling limit, and fixing the times removes it.
\pause
\item Scaling to real missions needs {\bf structure}: repeating sessions, invariants, and
      eventually assume-guarantee contracts per cycle.
\pause
\item The hard part is not the solver. It is knowing what the spec {\bf means}: \mexec{}'s semantics,
      and which properties are worth stating.
\end{enumerate}

\vspace{0.4cm}
{\centering\scriptsize Docs and source: {\tt https://nasa-jpl.github.io/tasksat/}\par}

\end{frame}

%----------------------------------------
\begin{frame}

\vfill
\centering {\Large\bf Questions?}

\vfill

\end{frame}
